\section{Introduction}

Consumer services increasingly depend on photographs as evidence: a guest reports a pest in a hotel room, a diner reports contamination in a meal, or a tenant reports damage in a short-term rental. These images are high-stakes enough to affect refunds, platform penalties, and insurance decisions, but they are often captured in uncontrolled settings and reviewed under time pressure. Modern image editing makes the threat sharper: an attacker does not need to synthesize a whole image. They can start from a real photograph and insert a small, plausible object in exactly the region that supports the claim.

This paper studies that localized threat. We focus on \emph{pixel-preserving local object insertion}: a real source image is edited only inside a small marked region, then the edited object is spliced back into the original frame. The surrounding image remains the genuine camera photograph. This is qualitatively different from whole-image AI generation and from traditional copy-move or splice attacks. A detector cannot rely on global diffusion artifacts when more than 99\% of the pixels may be real, and it cannot rely on a simple JPEG-versus-PNG shortcut when the forged and real images share the same source and post-processing.

Existing benchmarks leave this setting under-specified. Whole-image AI-generated image benchmarks such as GenImage evaluate whether a complete image was generated by a model \citep{zhu2023genimage}. Manipulation localization benchmarks and toolkits evaluate masks and image-level decisions across broader tampering categories \citep{guillaro2022trufor,ma2024imdldbenco}. Recent inpainting datasets move closer to our setting, but their domains and protocols are not built around consumer-claim adjudication \citep{chen2024gim,wang2025opensdi,yan2025cocoinpaint}. At the same time, bias analyses show that AI-image detectors often learn dataset construction artifacts rather than semantic forgery evidence \citep{grommelt2024fakejpeg}, and inpainting-exchange studies show that detectors can collapse when global generative traces are removed \citep{nebioglu2026inpaintingexchange}. CLAIMFORGE is designed to stress exactly this gap.

We make four contributions in this draft:
\begin{itemize}
    \item We define CLAIMFORGE, a benchmark task for localized AI object-insertion forgeries in consumer-claim imagery, with paired image-level detection and pixel-level localization targets.
    \item We specify a pixel-preserving construction pipeline that keeps the real source image as the reference context, changes only the inserted object region, and includes controls for file-format, object-presence, and paste-back artifacts.
    \item We propose an evaluation protocol that covers image manipulation localization models, whole-image AI-generated image detectors, inpainting-specific detectors, multimodal forensic reasoning models, human inspection, and provenance mechanisms.
    \item We report the current data-construction state of the repository: 600 screened candidate source images, 297 complete restaurant edit slots, 22 pilot HunyuanImage-3 generations, and insertion regions whose median area is 0.232\% of the source image.
\end{itemize}

This is an initial benchmark draft rather than a completed detector study. The paper therefore separates construction facts from planned evaluation results. We do not claim that existing detectors fail on CLAIMFORGE until the full baseline suite has been run under the protocol in Section~\ref{sec:protocol}. Instead, the draft records the task definition, the rationale for the stress test, the dataset-construction rules, and the experiments required to make the final claim defensible.
